\subsection{Compatibility, Not Optimizer Ranking}
\label{subsec:compatibility-not-ranking}

The central message of this paper is not that one optimizer is universally better
for uncertainty estimation. The relevant object is the compatibility between the
geometry induced by the optimizer and the geometry read by the detector. An
optimizer can produce a representation that remains favorable for logit-based
scores while changing feature-space properties that are important for distance,
density, neighborhood, angular, or subspace readouts. Conversely, a detector that
works well under one optimizer-induced geometry may become less reliable when the
same model family is trained with a different update rule.

This view changes how optimizer improvements should be evaluated. Accuracy,
calibration, or logit-level OOD behavior do not by themselves certify
feature-based uncertainty. When downstream reliability depends on penultimate
features, the optimizer should be evaluated together with the representation
geometry it induces and the detector family that will read that geometry. This is
especially important when comparing coupled and decoupled weight decay or
adaptive and non-adaptive optimizers, because these choices can alter class
structure, feature norms, covariance spectra, and local feature-bank geometry
without necessarily producing a large change in ID accuracy.

The compatibility view also avoids an overly broad conclusion. If a detector gap
appears under one optimizer, this does not imply that the optimizer is globally
bad or that the detector is intrinsically weak. It indicates that the detector's
readout channel and the learned geometry may be mismatched. The goal is therefore
to identify which geometry channel carries the mismatch and whether a diagnostic
control can localize it.

\subsection{What the Diagnostics Can and Cannot Identify}
\label{subsec:what-diagnostics-identify}

The diagnostic interventions in this work are not proposed as new OOD detectors
or universal fixes. Their role is to perturb or stabilize a specific geometry
channel so that detector behavior can be interpreted more mechanistically. L2
normalization removes radial feature-norm variation before distance-based
scoring. Covariance controls such as tied, diagonal, and shrinkage estimates
alter the conditioning and precision structure used by Mahalanobis-style and
DDU-style Gaussian feature-density scores. PCA and residual-subspace diagnostics
separate principal ID structure from directions that may contain nuisance or
shift-related variation.

A recovery after a diagnostic intervention narrows the set of plausible failure
channels, but it does not prove a complete causal mechanism by itself. For
example, if raw Mahalanobis or raw kNN improves after L2 normalization, radial or
norm-coupled geometry is implicated. However, the remaining behavior may still
depend on class separation, covariance anisotropy, local density, or OOD
severity. Similarly, if shrinkage improves a Gaussian density score, covariance
estimation or conditioning is involved; if the gap remains, covariance correction
alone is not sufficient to explain the mismatch.

The strongest evidence for the proposed framework would therefore come from an
aligned pattern across multiple measurements: optimizer choice changes the ID
geometry fingerprint, detector families respond according to their readout
channels, and diagnostic controls remove or reduce the corresponding failure
mode. A single geometry metric or a single detector ablation should not be
overinterpreted as a complete explanation.

\subsection{Limitations and Scope}
\label{subsec:limitations-scope}

This study is intentionally controlled. Its purpose is to test the
optimizer--geometry--detector pathway, not to provide an exhaustive benchmark
over all OOD detectors, architectures, datasets, or training recipes. The main
claims should therefore be limited to the evaluated optimizer families,
architectures, datasets, and detector implementations. Broader datasets,
larger-scale models, pretrained regimes, and additional OOD shifts are important
settings for testing how far the compatibility view generalizes.

The comparisons should also be interpreted as accuracy-controlled rather than
perfectly accuracy-matched. Residual differences in ID accuracy, calibration,
training stability, or hyperparameter sensitivity may contribute to detector
behavior and should be reported alongside geometry and OOD metrics. In
particular, if the AdamW-to-Adam interpolation changes ID accuracy strongly, it
should be interpreted as a joint optimization-and-geometry effect rather than a
clean geometry intervention.

Detector coverage is another limitation. DDU-style Gaussian feature-density
scores in this paper refer to the post-hoc Gaussian density scoring component,
not necessarily a faithful reproduction of the original DDU training recipe.
Likewise, CTM-, NECO-, prototype-, and residual-subspace diagnostics should be
treated as extended probes unless their score conventions, PCA dimensions,
class-selection rules, and implementation details are fixed within the same
evaluation protocol.

Finally, the geometry--detector link studied here is diagnostic rather than a
complete causal proof. Stronger causal claims would require direct interventions
that change one geometry channel while holding others fixed, tighter ID-accuracy
bands, and broader architecture--dataset coverage. The contribution of this
paper is to make the optimizer-induced geometry pathway explicit and to provide a
testable diagnostic structure for evaluating feature-based uncertainty.
